%\section{PAGINA 6}

%\section{Continuación del Resumen}
\begin{itemize}[series=resumenlist, resume]
    \item Este es un nuevo ítem que continúa la lista de la página anterior.
    \item Y aquí puedes agregar otro ítem más.
\end{itemize}

\section{Autoexamen}
\begin{itemize}
    \item Considerar el circuito de la figura 27-2a. $V=120 V$, $R_1=10$, 
    $R_2=390 \Omega$, $R_3=600 \Omega$, $R_L=270 \Omega$. Se supone que 
    la resistencia interna de la batería V es 0 en el circuito 
    equivalente de Norton.
    \begin{enumerate}
        \item $I_N = ..................A$.
        \item $R_N = .................. \Omega$.
        \item $I_L = ..................A$.
    \end{enumerate}
    \item Se considera el circuito de la figura 27-5a. $V_1=45 V$, $V_2=75 V$, y 
    se supone que la resistencia interna de estas baterías de cero, 
    $R_1=450 \Omega$, $R_2=1500 \Omega$, $R_3 (la carga) =470 \Omega$.
    En el circuito equivalente de Norton,
    \begin{enumerate}
        \item $I_N = ..................A$.
        \item $R_N = .................. \Omega$.
        \item $I_L = ...................A$.
    \end{enumerate}
\end{itemize}

\section{Materiales Necesarios}
\begin{itemize}
    \item Equipo: Dos fuentes de c.c. de tensión regulada.
    \item Dos miliamperímetros de varios alcances o dos VOM capaces de 
    medir hasta \SI{50}{\milli\ampere}.
    \item EVM.
    \item Resistencias: \SI{0.5}{\watt} \SI{3300}{\ohm}, \SI{0.5}{\watt} 
    \SI{5100}{\ohm}, \SI{10000}{\ohm} y otras que requieran la operación 8 (discrecional).
    \item Diversos: Potenciómetros 10000 $\Omega$ 2 W.
\end{itemize}

\section{Procedimiento}
\begin{enumerate}[series=procedimientolist]
    \item Conectar el circuito de la figura 27-6a. $S_1$ está abierto. 
    Conectar las fuentes de alimentación reguladas $V_1$ y $V_2$ y 
    ajustar $V_1=40 V$ y $V_2=20 V$. 
    \item $S_2$ está abierto. Cerrar $S_1$. Cerciorarse de que $V_1$ y 
    $V_2$ se matienen en $40 V$ y $20 V$ respectivamente. Medir con el 
    amperímetro M la corriente $I_L$ en $R_3$. Anotarla en la tabla 27-1, en 
    la columna \textbf{$I_L$ medida circuito original}.
    \item Cerrar $S_2$ cortocircuitando $R_3$. Medir la corriente $I_N$ que 
    suministra la fuente $V_1$ al cortocircuito. Anotarla en la tabla 27-1, en la 
    columna \textbf{$I_N$ medida circuito original}.
    \item Suprimir $V_1$ y $V_2$ y reemplazarlas por sus resistencias internas 
    (cortocircuitos). Eliminar $R_3$ en el circuito. Cortocicuitar los nodos 
    D y F entre si, así como los A y B. $S_1$ está aun cerrado, $S_2$ abierto. 
    Medir la resistencia entre los nodos F y G y anotarla en la tabla 27-1, en la 
    columna \textbf{$R_N$ medida circuito original}. Esta es la resistencia en shunt 
    con el generador equivalente de Norton.
    \item Conectar el circuito de la figura 27-6b. $S_1$ está abierto. 
    $R_3$ es la resistencia de carga de 5100 ohmios del circuito anterior. 
    $M_1$ y $M_2$ son miliamperímetros para medir la corriente de Norton 
    y la corriente de carga respectivamente. Ajustar la salida de tensión 
    V en cero voltios.
\end{enumerate}